%%%%%%%%%%%%%%%%%%%%%%%%%%%%%%%%%%%%%%%%%
% Cleese Assignment
% Structure Specification File
% Version 1.0 (27/5/2018)
%
% This template originates from:
% http://www.LaTeXTemplates.com
%
% Author:
% Vel (vel@LaTeXTemplates.com)
%
% License:
% CC BY-NC-SA 3.0 (http://creativecommons.org/licenses/by-nc-sa/3.0/)
% 
%%%%%%%%%%%%%%%%%%%%%%%%%%%%%%%%%%%%%%%%%

%----------------------------------------------------------------------------------------
%	PACKAGES AND OTHER DOCUMENT CONFIGURATIONS
%----------------------------------------------------------------------------------------

\usepackage{fontawesome5}

\usepackage{pifont}

\usepackage{lastpage} % Required to determine the last page number for the footer

\usepackage{enumitem}

\usepackage{amssymb}

\usepackage{verbatim}

\usepackage{mathrsfs}

\usepackage{float}

\usepackage{graphicx} % Required to insert images

\setlength\parindent{0pt} % Removes all indentation from paragraphs

\usepackage{tabularray}

\usepackage[most]{tcolorbox} % Required for boxes that split across pages

\usepackage{booktabs} % Required for better horizontal rules in tables

\usepackage{colortbl} % Required for \cellcolor in tables

\usepackage{siunitx}

\usepackage{codehigh}

\usepackage[normalem]{ulem}
\usepackage{cancel}

\usepackage{listings} % Required for insertion of code

\tcbset{
	hwrcodeboxstyle/.style={
		colback=gray!5,
		colframe=black!65,
		boxrule=0.45pt,
		arc=1.4mm,
		left=6pt,
		right=6pt,
		top=6pt,
		bottom=6pt
	}
}

\lstdefinelanguage{Rnote}{
	language=R,
	morekeywords={
		mutate,summarise,summarize,filter,group_by,ungroup,across,
		case_when,if_else,coalesce,replace_na,left_join,inner_join,
		full_join,semi_join,anti_join,pivot_longer,pivot_wider,
		ggplot,aes,labs,facet_wrap
	}
}

\lstdefinestyle{hwrcode}{
	language=Rnote,
	basicstyle=\ttfamily\scriptsize,
	keywordstyle=\color{blue!75!black},
	commentstyle=\itshape\color{teal!55!black},
	stringstyle=\color{red!70!black},
	showstringspaces=false,
	columns=fullflexible,
	keepspaces=true,
	tabsize=4,
	gobble=1,
	breaklines=true,
	breakatwhitespace=false,
	breakautoindent=true,
	breakindent=2em,
	aboveskip=0pt,
	belowskip=0pt,
	frame=none
}

\newtcblisting{rcodebox}{
	enhanced,
	breakable,
	listing only,
	hwrcodeboxstyle,
	listing options={style=hwrcode}
}

\usepackage{etoolbox} % Required for if statements

\usepackage{mathrsfs}

\usepackage{amssymb}

\usepackage{needspace}

\usepackage{capt-of}

\usepackage{soul}

\usepackage{pdfpages}

\usepackage{xcolor}

\usepackage{hyperref}

\usepackage{dsfont}

\usepackage{mathtools}

\usepackage{comment}

\UseTblrLibrary{booktabs}
\UseTblrLibrary{siunitx}
\providecommand{\tinytableTabularrayUnderline}[1]{\underline{#1}}
\providecommand{\tinytableTabularrayStrikeout}[1]{\sout{#1}}
\NewTableCommand{\tinytableDefineColor}[3]{\definecolor{#1}{#2}{#3}}

%----------------------------------------------------------------------------------------
%	MARGINS
%----------------------------------------------------------------------------------------
\synctex=1
\usepackage{geometry} % Required for adjusting page dimensions and margins

\geometry{
	paper=letterpaper, % Change to letterpaper for US letter
	top=1.2in, % Top margin
	bottom=1.2in, % Bottom margin
	left=1in, % Left margin
	right=1in, % Right margin
	headheight=14pt, % Header height
	footskip=1.4cm, % Space from the bottom margin to the baseline of the footer
	headsep=1.2cm, % Space from the top margin to the baseline of the header
	%showframe, % Uncomment to show how the type block is set on the page
}

%----------------------------------------------------------------------------------------
%	FONT
%----------------------------------------------------------------------------------------

\usepackage[T1]{fontenc}
\usepackage[utf8]{inputenc}
% Core math and theorem support.
\usepackage{amsmath, amssymb, amsthm, mathtools}

%----------------------------------------------------------------------------------------
%	HEADERS AND FOOTERS
%----------------------------------------------------------------------------------------

\usepackage{fancyhdr} % Required for customising headers and footers

\pagestyle{fancy} % Enable custom headers and footers

\lhead{\small\assignmentClass\ifdef{\assignmentClassInstructor}{\ (\assignmentClassInstructor):}{}\ \assignmentTitle} % Left header; output the instructor in brackets if one was set
\chead{} % Centre header
\rhead{\small\ifdef{\assignmentAuthorName}{\assignmentAuthorName}{\ifdef{\assignmentDueDate}{Due\ \assignmentDueDate}{}}} % Right header; output the author name if one was set, otherwise the due date if that was set

\lfoot{} % Left footer
\cfoot{\small Page\ \thepage\ of\ \pageref{LastPage}} % Centre footer
\rfoot{} % Right footer

\renewcommand\headrulewidth{0.5pt} % Thickness of the header rule
\providecommand{\subqnumbering}{\arabic} % Default style: 1,2,3,...
\providecommand{\subsubqnumbering}{\alph}
\renewcommand{\thefootnote}{\arabic{footnote}}

%----------------------------------------------------------------------------------------
%	MODIFY SECTION STYLES
%----------------------------------------------------------------------------------------

\usepackage{titlesec} % Required for modifying sections

%------------------------------------------------
% Section

\titleformat
{\section} % Section type being modified
[block] % Shape type, can be: hang, block, display, runin, leftmargin, rightmargin, drop, wrap, frame
{\Large\bfseries} % Format of the whole section
{\assignmentQuestionName~\thesection} % Format of the section label
{6pt} % Space between the title and label
{} % Code before the label

\titlespacing{\section}{0pt}{0.5\baselineskip}{0.5\baselineskip} % Spacing around section titles, the order is: left, before and after

%------------------------------------------------
% Subsection

\titleformat
{\subsection} % Section type being modified
[block] % Shape type, can be: hang, block, display, runin, leftmargin, rightmargin, drop, wrap, frame
{\itshape} % Format of the whole section
{(\subqnumbering{subsection})} % Format of the section label
{4pt} % Space between the title and label
{} % Code before the label

\titlespacing{\subsection}{0pt}{0.5\baselineskip}{0.5\baselineskip} % Spacing around section titles, the order is: left, before and after

\renewcommand\thesubsection{(\subqnumbering{subsection})}


% Subsubsection

\titleformat
{\subsubsection} % Section type being modified
[block] % Shape type
{\itshape} % Format of the whole section
{(\subsubqnumbering{subsubsection})} % Format of the section label
{4pt} % Space between the title and label
{} % Code before the label

\titlespacing{\subsubsection}{0pt}{0.5\baselineskip}{0.5\baselineskip}

\renewcommand\thesubsubsection{(\subsubqnumbering{subsubsection})}


%----------------------------------------------------------------------------------------
%	CUSTOM QUESTION COMMANDS/ENVIRONMENTS
%----------------------------------------------------------------------------------------
% In your preamble, define a counter for numbered questions:
\newcounter{question}
\newcounter{subquestion}

% Environment for numbered questions:
\newenvironment{question}{
  \renewcommand{\thefootnote}{\arabic{footnote}}
  \needspace{7\baselineskip}
  \setcounter{subsection}{0} % Reset subquestion counter for each new question
  \refstepcounter{question} % Increment the question counter
  \vspace{0.5\baselineskip} % Whitespace before the question
  \section*{\assignmentQuestionName\ \thequestion} % Use a starred section to avoid the section counter interfering
  \lfoot{} % Optional: adjust footer if necessary
}{
  \lfoot{}
}

% Environment for unnumbered questions:
\newenvironment{unnumberedquestion}{
  \renewcommand{\thefootnote}{\arabic{footnote}}
  \needspace{7\baselineskip}
  \vspace{0.5\baselineskip} % Whitespace before the question
  \setcounter{subsection}{0} % Reset subquestion counter for each new question
  \lfoot{}
}{
  \lfoot{}
}

%------------------------------------------------

% Environment for subquestions, takes 1 argument - the name of the section
\newenvironment{subquestion}[1]{
	\renewcommand{\thefootnote}{\arabic{footnote}}
    \needspace{7\baselineskip}
	\subsection{#1}
}{
}

%------------------------------------------------

% Environment for subsubquestions
\newenvironment{subsubquestion}[1]{
    \renewcommand{\thefootnote}{\arabic{footnote}}
    \needspace{5\baselineskip}
    \subsubsection{#1}
}{
}
%------------------------------------------------

% Command to print a question sentence
\newcommand{\questiontext}[1]{
	\renewcommand{\thefootnote}{\arabic{footnote}}
	{\itshape #1}
	\ifnum\lastnodetype=3 % Check if last node was display math
        \vspace{-2\baselineskip} % Reduce the extra space
    \fi
}

%------------------------------------------------

% Command to print a box that breaks across pages with the question answer
\newcommand{\answer}[1]{
	\renewcommand{\thefootnote}{\arabic{footnote}}
	\begin{tcolorbox}[
        breakable, 
        enhanced,
        parbox=false,
		before upper={\setlength{\parskip}{0.5\baselineskip}},
		after skip=1.0\baselineskip
    ]
		\renewcommand{\thefootnote}{\arabic{footnote}}
        #1
    \end{tcolorbox}
}

%------------------------------------------------

% Command to print a box that breaks across pages with the space for a student to answer
\newcommand{\answerbox}[1]{
	\begin{tcolorbox}[breakable, enhanced]
		\renewcommand{\thefootnote}{\arabic{footnote}}
		\vphantom{L}\vspace{\numexpr #1-1\relax\baselineskip} % \vphantom{L} to provide a typesetting strut with a height for the line, \numexpr to subtract user input by 1 to make it 0-based as this command is
	\end{tcolorbox}
}

%------------------------------------------------

% Command to print an assignment section title to split an assignment into major parts
\newcommand{\assignmentSection}[1]{
	{
		\centering % Centre the section title
		\vspace{2\baselineskip} % Whitespace before the entire section title
		
		
		\Large{\textbf{#1}} % Section title, forced to be uppercase
		
		
		\vspace{\baselineskip} % Whitespace after the entire section title

	}
}


%----------------------------------------------------------------------------------------
%	TITLE PAGE
%----------------------------------------------------------------------------------------

\author{\textbf{\assignmentAuthorName}} % Set the default title page author field
\date{} % Don't use the default title page date field

\title{
	\thispagestyle{empty} % Suppress headers and footers
	\vspace{0.2\textheight} % Whitespace before the title
	\textbf{\assignmentClass:\ \assignmentTitle}\\[-4pt]
	\ifdef{\assignmentDueDate}{{\small Due\ on\ \assignmentDueDate}\\}{} % If a due date is supplied, output it
	\ifdef{\assignmentClassInstructor}{{\large \textit{\assignmentClassInstructor}}}{} % If an instructor is supplied, output it
	\vspace{0.32\textheight} % Whitespace before the author name
}


\newcommand{\ind}{\perp\!\!\!\perp} 
\newcommand*\N{\mathbb{N}}
\newcommand*\R{\mathbb{R}}
\newcommand*\Q{\mathbb{Q}}
\newcommand*\Z{\mathbb{Z}}
\newcommand*\E{\mathbb{E}}
\newcommand*\oldL{\L}
\renewcommand*\L{\mathbb{L}}
\newcommand*\oldP{\P}
\renewcommand*\P{\mathbb{P}}
\let\O\varnothing
\let\oldemptyset\emptyset
\let\emptyset\varnothing

\DeclareMathOperator*{\argmax}{arg\,max} % thin space, limits underneath in displays
\DeclareMathOperator*{\argmin}{arg\,min} % thin space, limits underneath in displays
\DeclareMathOperator{\opvec}{vec}
\DeclareMathOperator*{\MSE}{MSE}
\DeclareMathOperator*{\Var}{Var}
\DeclareMathOperator*{\Cov}{Cov}
\DeclareMathOperator*{\Bias}{Bias}
\DeclareMathOperator*{\Corr}{Corr}
\DeclareMathOperator*{\plim}{plim}
\DeclareMathOperator*{\esssup}{ess\,sup}
\DeclareMathOperator*{\essinf}{ess\,inf}
